\section{Tổng quan và đặc tả dự án}
\OwnerNote{Trần Nhật Bảo Anh -- 2410155 (toàn bộ Phần I; kiểm soát nội dung toàn báo cáo)}

\subsection{Giới thiệu đề tài}

\Placeholder{Giới thiệu ngắn gọn Smart E-Mobility Hub, lý do nhóm thực hiện tài liệu này và những nội dung được xác định trong Submission 1.}

\subsection{Bối cảnh và vấn đề cần giải quyết}

\subsubsection*{Bối cảnh dự án}

\Placeholder{Viết bối cảnh di chuyển trong Khu đô thị ĐHQG-HCM, vai trò của Metro, các trường đại học, ký túc xá, khu dịch vụ và mạng lưới Mobility Hub. Không đưa framework hoặc giải pháp cài đặt vào phần này.}

\subsubsection*{Các vấn đề hiện tại}

\noindent\textbf{PB-01 -- \Placeholder{Tên vấn đề}.}
\Placeholder{Mô tả biểu hiện của vấn đề, stakeholder bị ảnh hưởng, nguyên nhân và hậu quả đối với việc di chuyển hoặc vận hành.}

\medskip
\noindent\textbf{PB-02 -- \Placeholder{Tên vấn đề}.}
\Placeholder{Mô tả biểu hiện của vấn đề, stakeholder bị ảnh hưởng, nguyên nhân và hậu quả đối với việc di chuyển hoặc vận hành.}

\subsubsection*{Lý do cần xây dựng hệ thống}

\Placeholder{Giải thích vì sao các vấn đề trên cần một hệ thống quản lý và điều phối thống nhất.}

\subsection{Stakeholder và nhu cầu}

\subsubsection*{Danh sách stakeholder}

\begin{table}[htbp]
  \centering
  \caption{Stakeholder của dự án}
  \label{tab:stakeholders}
  \small
  \begin{tabularx}{\textwidth}{|L{20mm}|L{34mm}|Y|L{28mm}|}
    \hline
    \rowcolor{gray!15}
    \textbf{Mã} & \centering\arraybackslash\textbf{Stakeholder} & \centering\arraybackslash\textbf{Vai trò và mối quan tâm} & \centering\arraybackslash\textbf{Actor tương ứng}\\
    \hline
    SH-STU & Sinh viên & \Placeholder{Cần điền} & ACT-STU\\
    \hline
    SH-OPR & Đơn vị vận hành & \Placeholder{Cần điền} & ACT-OPR\\
    \hline
    SH-MGT & Đơn vị quản lý & \Placeholder{Cần điền} & \Placeholder{Có/Không}\\
    \hline
    SH-MNT & Nhân viên kỹ thuật & \Placeholder{Cần điền} & ACT-MNT\\
    \hline
  \end{tabularx}
\end{table}

\subsubsection*{Stakeholder needs}

\begin{longtable}{|L{23mm}|L{28mm}|L{32mm}|L{23mm}|L{23mm}|}
  \caption{Nhu cầu của stakeholder}\label{tab:stakeholder-needs}\\
  \hline
  \rowcolor{gray!15}
  \textbf{Need ID} & \textbf{Stakeholder} & \centering\arraybackslash\textbf{Nhu cầu hoặc mong đợi} & \textbf{Problem} & \textbf{Objective}\\
  \hline
  \endfirsthead
  \hline
  \rowcolor{gray!15}
  \textbf{Need ID} & \textbf{Stakeholder} & \centering\arraybackslash\textbf{Nhu cầu hoặc mong đợi} & \textbf{Problem} & \textbf{Objective}\\
  \hline
  \endhead
  SN-STU-01 & SH-STU & \Placeholder{Cần viết} & PB-XX & OBJ-XX\\
  \hline
  SN-OPR-01 & SH-OPR & \Placeholder{Cần viết} & PB-XX & OBJ-XX\\
  \hline
\end{longtable}

\subsection{User Stories}

\begin{longtable}{|L{23mm}|L{67mm}|L{25mm}|L{22mm}|}
  \caption{User Stories của dự án}\label{tab:user-stories}\\
  \hline
  \rowcolor{gray!15}
  \textbf{User Story ID} & \centering\arraybackslash\textbf{Nội dung} & \textbf{Need} & \textbf{Phân hệ}\\
  \hline
  \endfirsthead
  \hline
  \rowcolor{gray!15}
  \textbf{User Story ID} & \centering\arraybackslash\textbf{Nội dung} & \textbf{Need} & \textbf{Phân hệ}\\
  \hline
  \endhead
  US-HUB-01 & Là một \Placeholder{actor}, tôi muốn \Placeholder{mục tiêu} để \Placeholder{giá trị nhận được}. & SN-XX-XX & SUB-HUB\\
  \hline
  US-OPS-01 & Là một \Placeholder{actor}, tôi muốn \Placeholder{mục tiêu} để \Placeholder{giá trị nhận được}. & SN-XX-XX & SUB-OPS\\
  \hline
\end{longtable}

\subsection{Mục tiêu dự án}

\noindent\textbf{OBJ-01 -- \Placeholder{Tên mục tiêu}.}
\Placeholder{Trình bày kết quả dự án muốn đạt, vấn đề được giải quyết và dấu hiệu dùng để đánh giá mục tiêu.}

\medskip
\noindent\textbf{OBJ-02 -- \Placeholder{Tên mục tiêu}.}
\Placeholder{Trình bày kết quả dự án muốn đạt, vấn đề được giải quyết và dấu hiệu dùng để đánh giá mục tiêu.}

\subsection{Phạm vi dự án}

\subsubsection*{Trong phạm vi}

\begin{itemize}
  \item \textbf{SCP-IN-01:} \Placeholder{Mô tả năng lực nghiệp vụ thuộc phạm vi và phân hệ chịu trách nhiệm.}
  \item \textbf{SCP-IN-02:} \Placeholder{Bổ sung các năng lực còn lại theo bảy phân hệ.}
\end{itemize}

\subsubsection*{Ngoài phạm vi}

\begin{itemize}
  \item \textbf{SCP-OUT-01:} \Placeholder{Mô tả nội dung không thuộc phạm vi và giải thích lý do loại trừ.}
\end{itemize}

\subsubsection*{Phạm vi phân tích và phạm vi MVP}

\Placeholder{Phân biệt toàn bộ nghiệp vụ được phân tích với các luồng được chọn để hiện thực trong MVP.}

\subsection{System Boundary}

\begin{itemize}
  \item \textbf{Phần mềm Smart E-Mobility Hub chịu trách nhiệm:} \Placeholder{Các quyết định, dữ liệu và chức năng nằm trong hệ thống.}
  \item \textbf{Con người chịu trách nhiệm:} \Placeholder{Các hành động vật lý hoặc quyết định không được tự động hóa.}
  \item \textbf{Hệ thống bên ngoài cung cấp:} \Placeholder{Dữ liệu và dịch vụ được xem là đầu vào, không thuộc quyền kiểm soát của nhóm.}
\end{itemize}

\DiagramPlaceholder{Context diagram và ranh giới hệ thống}{fig:context-diagram}

\clearpage
\subsection{Phân rã hệ thống thành bảy phân hệ nghiệp vụ}

Trong báo cáo này, tiền tố \texttt{SUB} là viết tắt của \textit{subsystem}. Mỗi mã \texttt{SUB-*} đại diện cho một phân hệ nghiệp vụ có phạm vi và người chịu trách nhiệm riêng. Sau phần giới thiệu này, báo cáo luôn ghi mã cùng tên đầy đủ ở tiêu đề; trong nội dung chi tiết có thể dùng mã ngắn để tránh lặp lại. Trên use-case diagram, system boundary phải ghi cả mã và tên phân hệ, không chỉ ghi một mã \texttt{SUB-*} đứng riêng.

\medskip
\noindent\textbf{SUB-HUB -- Tra cứu Hub và tài nguyên.}\par
\Placeholder{Mô tả trách nhiệm, dữ liệu chính và ranh giới của phân hệ tra cứu Hub.}

\medskip
\noindent\textbf{SUB-RES -- Đặt phương tiện dùng chung.}\par
\Placeholder{Mô tả trách nhiệm, dữ liệu chính và ranh giới của phân hệ đặt phương tiện.}

\medskip
\noindent\textbf{SUB-TRIP -- Quản lý chuyến đi.}\par
\Placeholder{Mô tả trách nhiệm, dữ liệu chính và ranh giới của phân hệ quản lý chuyến đi.}

\medskip
\noindent\textbf{SUB-PARK -- Quản lý chỗ đỗ xe cá nhân.}\par
\Placeholder{Mô tả trách nhiệm, dữ liệu chính và ranh giới của phân hệ chỗ đỗ.}

\medskip
\noindent\textbf{SUB-CHG -- Quản lý sạc và lịch sạc.}\par
\Placeholder{Mô tả trách nhiệm, dữ liệu chính và ranh giới của phân hệ sạc.}

\medskip
\noindent\textbf{SUB-OPS -- Giám sát và điều phối vận hành.}\par
\Placeholder{Mô tả trách nhiệm, dữ liệu chính và ranh giới của phân hệ vận hành.}

\medskip
\noindent\textbf{SUB-SIM -- What-if Simulation và khuyến nghị.}\par
\Placeholder{Mô tả trách nhiệm, dữ liệu chính và ranh giới của phân hệ mô phỏng.}

\subsection{Assumptions và Constraints}

\subsubsection*{Assumptions}

\noindent\textbf{ASM-GEN-01 -- \Placeholder{Tên giả định}.}
\Placeholder{Nêu điều nhóm đang tạm xem là đúng trong quá trình phân tích. Nếu giả định không còn đúng, chỉ rõ FR, BR hoặc use case cần xem lại.}

\medskip
\noindent\textbf{ASM-GEN-02 -- \Placeholder{Tên giả định}.}
\Placeholder{Bổ sung các giả định về dữ liệu, hạ tầng, người dùng hoặc cách vận hành có ảnh hưởng trực tiếp đến phạm vi hệ thống.}

\subsubsection*{Constraints}

\noindent\textbf{CON-GEN-01 -- \Placeholder{Tên ràng buộc}.}
\Placeholder{Mô tả giới hạn bắt buộc xuất phát từ đề bài, thời gian, nguồn lực, dữ liệu hoặc công nghệ; ghi rõ nguồn của ràng buộc.}

\medskip
\noindent\textbf{CON-GEN-02 -- \Placeholder{Tên ràng buộc}.}
\Placeholder{Bổ sung ràng buộc khác và ảnh hưởng của nó đến cách nhóm phân tích hoặc hiện thực hệ thống.}

\subsection{Thuật ngữ và trạng thái dùng chung}

\Placeholder{Chèn bảng rút gọn từ project-conventions.md, chỉ giữ các thuật ngữ thực sự được dùng trong báo cáo.}

\clearpage
